\begin{lemma}[Situation 7.2, Subcase 1.2]
Every \noderef{HierarchySituationTwoCaseOneTwoData} is impossible.
Equivalently, Situation 7.2, Case 1 cannot occur when the last-edge color
\(\theta=\varphi_n(e_p)\) is missing at \(y_{p-1}\).
\end{lemma}
\begin{proof}
Write \(H=C.\mathit{hierarchy}\), let \(q\) be its last stored index, and
put
\[
 A=\noderef{HierarchyStarredStage}(H,q),\qquad
 \ell=\noderef{HierarchyFullTailLength}(C),
\]
\[
 B=\noderef{HierarchyFullTailPrefix}(C,\ell-2),\qquad
 B^+=\noderef{HierarchyFullTailPrefix}(C,\ell-1),\qquad
 D_{n,q}=\noderef{HierarchyStoredColorsAt}(H,q).              \tag{1}
\]
The exact position fields in the supplied datum, its inequality
\(\ell\ge2\), and the definitions in (1) give entrywise
\noderef{TreeSequenceInitialSegment}s
\[
 A\preccurlyeq B\prec B^+\prec T,\qquad
 |B|=|T|-2,\quad |B^+|=|T|-1.                               \tag{2}
\]
Thus the last vertices of \(B^+\) and \(T\) are respectively
\(y_{p-1}\) and \(y_p\), and the last individual edge of \(T\) is
\(e_p\).  Apply
\noderef{HierarchyCounterexamplePenultimateElementary} to the properness,
exact one-uncoloredness, complete lower-rung package, fewer-divider
induction, lexicographic minimality, and large-set fields of the supplied
\noderef{HierarchyMinimumCounterexampleContext}.  Its exact conclusion is
\[
             V(B^+)\text{ is elementary under }\varphi_n.    \tag{3}
\]

We first derive all parts of manuscript (1) and (7).  The datum says that
\(\alpha\in D_{n,q}\) and that \(\alpha\) is missing at
\(y_{p-1}\).  Expanding \noderef{HierarchyStoredColorsAt} shows that
\(\alpha\) is not missing on \(A\).  By (3), no vertex of \(B\) misses
\(\alpha\); hence the totalized
\noderef{HierarchyFirstMissingPosition} of \(\alpha\) is exactly the
position of \(y_{p-1}\).  The equation
\(C.\mathit{good}.\Gamma^q_\alpha=\{\gamma_1,\gamma_2\}\) and the
avoidance field of \noderef{GoodHierarchy} consequently give, for
\(j=1,2\),
\[
 \gamma_j\in\overline\varphi_n(A),\qquad
 \gamma_j\notin\varphi_n\langle B^+-A\rangle .               \tag{4}
\]
The two colors are distinct, as recorded in the datum.  Applying
\noderef{HierarchyTailMissingColorExclusion} to the first
\(\alpha\)-missing vertex \(y_{p-1}\) also shows that \(\alpha\) colors
no exact edge of \(B^+-A\), and no later exact edge of \(T-A\), apart
from the possible first edge when \(q=0\), \(A=T_n\), and that edge is
the individual SE/RE seed \(f_n\) carrying \(\alpha\in D_n\).  The last
edge is not exceptional: it is incident with \(y_p\), while \(\alpha\)
is missing there.

The Subcase 1.2 field says that \(\theta\) is missing at
\(y_{p-1}\).  Elementarity (3) therefore gives
\[
                  \theta\notin\overline\varphi_n(B).         \tag{5}
\]
Every color in a last Gamma pair is missing on the corresponding stage
and hence has a witness in \(A\).  Such a witness is distinct from
\(y_{p-1}\) by (2).  Thus (3) and the missingness of \(\theta\) at
\(y_{p-1}\) imply
\[
 \theta\notin\Gamma^q,\qquad
 \theta\ne\gamma_1,\gamma_2.                                \tag{6}
\]
Apply \noderef{HierarchyTailMissingColorExclusion} at the exact last
vertex of \(B^+\).  Together with (5), it proves that \(\theta\) colors
no individual edge of \(B^+-A\), except possibly the first edge in the
divider-zero SE/RE case, where that edge is exactly \(f_n\) and
\(\theta\in D_n\).  Equations (5)--(6) and this exact seed alternative
are the full formal content of manuscript equation (7).

We now change \(\alpha\) and \(\gamma_1\).  Both are missing on
\(B^+\): \(\gamma_1\) has a witness in \(A\) by (4), and \(\alpha\)
is missing at \(y_{p-1}\).  They are distinct because \(\alpha\) is not
missing on \(A\).  Apply \noderef{ChangeLemma} with the ordered colors
\((\gamma_1,\alpha)\).  Its conditional order premise is active:
the first \(\gamma_1\)-missing vertex lies in \(A\), the first
\(\alpha\)-missing vertex is \(y_{p-1}\), and (4) supplies the exact
unused interval.  For \(q>0\), \(\gamma_1\) is missing on
\(T_{n,q}=A\); for \(q=0\), it is missing on \(T_n\).  This is exactly
the Cases 1/2 proposition required by both the change lemma and
\noderef{StableChangeCasesOneTwo}; no stored-set nonmembership is being
used as stage missingness.  With (3) and the six minimum-context inputs,
the change lemma gives
\[
 K_{\alpha,\gamma_1}(v_{\gamma_1})
       =K_{\alpha,\gamma_1}(y_{p-1}),                         \tag{7}
\]
where \(v_{\gamma_1}\in V(A)\) is the first
\(\gamma_1\)-missing vertex.

Let \(Q_1=K_{\alpha,\gamma_1}(y_p)\).  The component (7) and \(Q_1\)
are disjoint.  If \(v_{\gamma_1}=y_{p-1}\), that common vertex misses
both affected colors, so expanding \noderef{MissingColors} and
\noderef{TwoColorIncidenceGraph} makes its component an isolated vertex;
it cannot contain the distinct \(y_p\).  Otherwise, equality of the two
components would put the three distinct vertices
\(v_{\gamma_1},y_{p-1},y_p\), missing respectively
\(\gamma_1,\alpha,\alpha\), into one component.  Properness and
\noderef{TwoColorIncidenceDegreeAtMostTwo} make each a degree-at-most-one
vertex, while \noderef{TwoColorComponentSpanningPath} gives a finite
spanning path with only two endpoints.  This is impossible.  The
dichotomy in \noderef{TwoColorComponentsDisjoint} gives the asserted
disjointness, and the spanning-path theorem also supplies the possibly
singleton \noderef{TwoColorPathComponent} certificate for \(Q_1\).

Define
\[
 \mu_1=\noderef{TwoColorComponentSwap}
       (G,\varphi_n,\gamma_1,\alpha,\operatorname{inl}y_p).   \tag{8}
\]
Apply \noderef{StableChangeCasesOneTwo} immediately, using the exact
stage-missing Cases 1/2 proposition for \(\gamma_1\), (3), (7)--(8),
the component separation, and all six context fields.  Retain its
\noderef{HierarchyWeaklyStable} and \noderef{ModColoring} conclusions.

We next verify the stronger hypotheses of
\noderef{StableChangeCasesOneTwoFullTail}.  By (4), \(\gamma_1\) is
absent from every tail edge through \(B^+\), and it is absent from
\(e_p\) by (6) and \(\varphi_n(e_p)=\theta\).  The color \(\alpha\) is
absent from every exact edge of \(T-A\), except for the individual
divider-zero seed described above.  In that exception the already
obtained weak stability contains a \noderef{StableColoring} on \(T_n\);
the seed is a \noderef{TreeIncidentEdge} carrying the protected stored
color \(\alpha\), so its incident-edge clause fixes that exact edge and
shows that it lies outside \(Q_1\).  Hence
\[
 \mu_1(T.\mathit{edge}(r))=\varphi_n(T.\mathit{edge}(r))
                  \quad(|A|\le r<|T|).                       \tag{9}
\]
For every vertex of \(B^+\), missing sets also agree.  Outside \(Q_1\)
no incident affected edge changes.  A vertex of \(Q_1\cap V(B^+)\)
cannot be an affected-color endpoint: by elementarity (3), any such
endpoint is the corresponding unique witness in the disjoint component
(7).  It is therefore internal, with one incident edge of each affected
color by properness, and the swap exchanges those two edges.  Thus
\[
 \overline\mu_1(u)=\overline\varphi_n(u)
                         \quad(u\in V(B^+)).                  \tag{10}
\]
For a stored color, \(\gamma_1\) cannot occur because it is missing on
\(A\).  The stored affected color \(\alpha\) has first witness
\(y_{p-1}\), and (10) preserves all candidates through that witness.
Every other color is unaffected everywhere.  Expanding
\noderef{HierarchyFirstMissingPosition}, including its fallback case,
therefore gives equality of the first positions of every color in
\(D_{n,q}\).  Equations (9)--(10), this equality, the path certificate,
and the exact Cases 1/2 premise verify every input of the full-tail
provider.

Consequently \noderef{StableChangeCasesOneTwoFullTail} returns an actual
\noderef{CorrespondingExtendedTashkinovTree} on the same exact \(T\),
maximum property for its replacement ETT, a replacement base
\(\mathcal B_1\), a hierarchy \(H_1\), a \noderef{GoodHierarchy}
witness \(g_1\), and
\[
             \noderef{SameGoodHierarchy}(C.\mathit{good},g_1). \tag{11}
\]
The swap makes \(\gamma_1\) missing at \(y_p\), including when \(Q_1\) is
the permitted singleton component.  Indeed, first suppose that an edge
\(f\) colored \(\gamma_1\) is incident with \(y_p\).  Properness makes
\(f\) the unique such edge.  Its edge-node is adjacent to the root
\(y_p\) in \noderef{TwoColorIncidenceGraph}, so it belongs to \(Q_1\),
and (8), unfolded using \noderef{TwoColorComponentSwap}, changes its color
from \(\gamma_1\) to \(\alpha\).  Every old incident
\(\gamma_1\)-edge is therefore removed, while no new incident
\(\gamma_1\)-edge can be created: such a new edge would have had old color
\(\alpha\), contrary to \(\alpha\in\overline\varphi_n(y_p)\) and the
definition of \noderef{MissingColors}.  Hence \(\gamma_1\) is missing at
\(y_p\) under \(\mu_1\).  In the complementary case no
\(\gamma_1\)-edge is incident with \(y_p\), so \(\gamma_1\) was already
missing there.  Since \(\alpha\) is missing there as well, the root has no
incident affected edge in the two-color incidence graph.  Thus \(Q_1\) is
the isolated singleton allowed by \noderef{TwoColorPathComponent}; the
swap changes no edge-node and preserves the \(\gamma_1\)-missingness of
\(y_p\).

The disjoint component (7) preserves the old \(\gamma_1\)-missing witness
in \(A\).  By (2), that witness and \(y_p\) are distinct, so these two
witnesses make \(T\) non-elementary.  Together with the returned
ETT, maximum property, base, hierarchy, and goodness, they give an actual
\noderef{HierarchyCounterexample} \(C_1\).

Package these data as a
\noderef{HierarchyReplacementPresentation} \(P_1\), constructing its
minimum context field by field.  Properness is
\noderef{TwoColorComponentSwapProper}; the pointwise swap preserves
`none', hence preserves \noderef{IsOneEdgeUncolored}.  Rung agreement in
the correspondence transports the complete universally quantified
lower-rung six-pack.  Equation (11) transports divider count, so the old
fewer-divider induction applies.  The terminal ordered tree is literally
\(T\), and (11) identifies its stages and divider positions; unfolding
\noderef{HierarchyCounterexampleMeasure} therefore identifies both
measure coordinates, so old lexicographic minimality applies.  Finally,
the terminal positive-slot stability in the correspondence, unfolded via
\noderef{StableColoring}, identifies missing sets on \(T_n\) and transports
the \(2n+11\) estimate.  These are every field of the replacement minimum
context.

The first swap also proves manuscript (8):
\[
                 \gamma_1\notin\mu_1\langle T-A\rangle .     \tag{12}
\]
Indeed, (4) and (6) excluded \(\gamma_1\) from every old tail edge, and
the only way for the swap to introduce \(\gamma_1\) is from an old
\(\alpha\)-edge in \(Q_1\).  There is no such tail edge except the
possible seed, and (9) fixed that individual edge.

We now change \(\gamma_1\) and \(\theta\) relative to \(P_1\).
Equation (10) preserves the \(\theta\)-missing vertex \(y_{p-1}\), while
\(\gamma_1\) remains missing at its witness in \(A\).  They are distinct
by (6).  Apply \noderef{HierarchyCounterexamplePenultimateElementary}
again to \(P_1.\mathit{minimumContext}\); it gives elementarity of the
same exact \(B^+\) under \(\mu_1\).  Apply \noderef{ChangeLemma} with
ordered colors \((\gamma_1,\theta)\).  The first color is missing on the
last ordinary stage (or on \(T_n\) when \(q=0\)), so its exact Cases 1/2
premise holds.  Its first witness lies in \(A\), theta's first witness is
\(y_{p-1}\) by (5), (10), and elementarity, and (12) supplies the required
unused interval.  The conclusion is
\[
 K^{\mu_1}_{\gamma_1,\theta}(v_{\gamma_1})
       =K^{\mu_1}_{\gamma_1,\theta}(y_{p-1}).                 \tag{13}
\]
The component in (13) is disjoint from the component rooted at \(y_p\).
The coincident-witness case is again an isolated vertex; otherwise the
three distinct missing endpoints
\(v_{\gamma_1},y_{p-1},y_p\), missing respectively
\(\gamma_1,\theta,\gamma_1\), contradict the finite spanning-path and
degree bounds.  Define
\[
 \mu_2=\noderef{TwoColorComponentSwap}
       (G,\mu_1,\gamma_1,\theta,\operatorname{inl}y_p).       \tag{14}
\]
Invoke \noderef{StableChangeCasesOneTwo} relative to \(P_1\) before using
any seed protection, with \(\gamma_1\) in the stage-missing slot, (13),
the component separation, and the six fields of the new context.  Retain
its weak stability and mod-coloring conclusions.

Only the prefix \(B^+\), not the final edge, is fixed by this swap.
Equation (12) excludes \(\gamma_1\) from all its exact tail edges, and
the theta-exclusion following (6) excludes \(\theta\) from
\(B^+-A\), apart from the possible exact divider-zero seed.  The retained
weak stability fixes that seed exactly as above.  Hence all individual
edges from \(A\) through \(B^+\) retain their values.  Repeating the
internal-vertex argument using (13) gives
\[
 \overline\mu_2(u)=\overline\mu_1(u)
                         \quad(u\in V(B^+)).                  \tag{15}
\]

We verify the route input of
\noderef{StableChangeCasesOneTwoSplitTail} with \(X=B^+\).  Inspect the
actual final certificate in the \noderef{TashkinovSeries} of \(P_1\)'s
replacement ETT.  The PE constructor supplies a
\(\mu_1\)-\noderef{TashkinovClosure} from its terminal ladder tree; the
SE/RE constructors supply the exact singleton connecting edge, literal
\noderef{TreeSequenceAdjoinEdge} seed, and closure from that seed.
The base and hierarchy prefix chain place the eligible source below
\(A\preccurlyeq B^+\).  Apply
\noderef{TashkinovAugmentationSequenceRestrict} to obtain the old route to
\(B^+\).  Below \(A\), use the returned
\noderef{HierarchyWeaklyStable}: its positive-divider
\noderef{TreeColoringInvariant}, or at divider zero its PE combined
stability together with \noderef{CombinedStableClosurePreserved} and the
entry decomposition of \noderef{OrderedTreeUnion}, transports each exact
step.  In the non-PE zero case \(A=T_n\), so only the protected exact seed
remains.  Above \(A\), fixed edge values and (15) let
\noderef{TashkinovAugmentationSequenceColoringTransport} carry the route
to \(\mu_2\).

The old suffix \(B^+\to T\) is the one
\noderef{TashkinovAugmentation} adding \((e_p,y_p)\) under \(\mu_1\):
the first swap fixed \(e_p\), its color is \(\theta\), and \(\theta\)
is missing at \(y_{p-1}\).  Equation (15) leaves both affected colors
missing on the split prefix.  For the stored-position premise,
\(\gamma_1\) is not stored because it is stage-missing; if \(\theta\)
is stored, (15) preserves its first witness \(y_{p-1}\), and every other
stored color is unaffected.  For the Gamma cutoff, (6) places \(\theta\)
in no pair, while any pair containing \(\gamma_1\) is the pair of
\(\alpha\) by pairwise disjointness in \noderef{GoodHierarchy}; the first
\(\alpha\)-missing position is \(y_{p-1}\), exactly \(|B^+|\).  Thus all
six transport inputs of the split-tail provider hold.

The split-tail theorem returns a second corresponding full-tree ETT,
maximum property, base \(\mathcal B_2\), good hierarchy \(g_2\), and
same-good relation relative to \(P_1\).  Compose this correspondence back
to the original ETT field by field.  Rung, ladder-tree, raw stored-set,
connecting-edge, kind, and generating-coloring equalities compose
directly.  At every positive coloring slot, use transitivity
\noderef{sc2} of \noderef{StableColoring}, after rewriting the protected
stored set by the intervening ladder-tree and raw-stored equalities.
Compose the two \noderef{SameGoodHierarchy} relations pointwise: stages
and divider positions, computed stored sets, and every indexed Gamma pair.
This produces an original-relative replacement presentation \(P_2\).
Its counterexample uses the two \(\theta\)-missing vertices
\(y_{p-1},y_p\): the swap at \(y_p\) changes the incident last edge from
\(\theta\) to \(\gamma_1\), while (15) preserves the first witness.
Its minimum context is constructed by the same five fieldwise transports
used for \(P_1\), now with the composed correspondence and same-good
relations.

Finally,
\[
 \mu_2(e_p)=\gamma_1,\qquad
 \theta\in\overline\mu_2(y_{p-1})\cap
             \overline\mu_2(y_p),\qquad
 \gamma_1\notin\overline\mu_2(y_{p-1}).                      \tag{16}
\]
The last assertion holds because \(e_p\) is not incident with
\(y_{p-1}\), (13) puts that vertex in the component disjoint from the
one swapped at \(y_p\), and it was an endpoint missing \(\theta\), hence
present in \(\gamma_1\), before the swap.  Split according to whether
\(\theta\) belongs to the last stored set of \(P_2\).  If it does, the
two-color cardinality field of \(g_2\) supplies distinct
\(\delta_1,\delta_2\) whose pair is
\(\Gamma^q_\theta\).  The common-tail fields transported by the composed
correspondence, (16), storedness, and this pair form an actual
\noderef{HierarchySituationTwoCaseOneData};
\noderef{HierarchySituationTwoSubcaseOneOne} gives a contradiction.  If
\(\theta\) is not stored, the same common-tail fields and (16) form an
actual \noderef{HierarchySituationTwoCaseTwoData}, which is impossible by
\noderef{HierarchySituationTwoSubcaseTwoOne}.  The storedness split is
exhaustive, completing Subcase 1.2.
\end{proof}
